\section{Hill Cipher}

\subsection{Konsep}

Cipher Hill menggunakan aljabar linear: plaintext dibagi menjadi blok berukuran $n$, lalu dikalikan dengan matriks kunci $K$ berukuran $n \times n$ secara modulo 26 \cite{stallings2017,trappe2006}. 

\[
\mathbf{C} = K \mathbf{P} \bmod 26
\]

Dekripsi membutuhkan invers matriks $K^{-1}$ modulo 26. Matriks $K$ harus invertible (determinannya relatif prima dengan 26).

\subsection{Matriks Kunci dan Invers}

Matriks kunci Hill $K$ berukuran $n \times n$ dengan elemen $k_{ij}$:
\[
K = \begin{bmatrix}k_{11} & k_{12} & \cdots & k_{1n} \\ k_{21} & k_{22} & \cdots & k_{2n} \\ \vdots & \vdots & \vdots \\ k_{n1} & k_{n2} & \cdots & k_{nn} \end{bmatrix}
\]

Invers matriks $K^{-1}$ dihitung dengan:
\[
K^{-1} = \text{adj}(K) \cdot \det(K)^{-1} \bmod 26
\]
dimana $\text{adj}(K)$ adalah adjoin dari kofaktor kofaktor elemen minor.

\subsection{Contoh 2x2}

Gunakan $K=\begin{bmatrix}3 & 3\\ 2 & 1\end{bmatrix}$. Plaintext HI:
\begin{itemize}
  \item Vektor: $(7, 8)$
  \item Enkripsi: $\begin{bmatrix}3 & 3\\ 2 & 1\end{bmatrix} \begin{bmatrix}7\\ 8\end{bmatrix} = \begin{bmatrix}7\\ 20\end{bmatrix} \equiv \begin{bmatrix}7\\ 20\end{bmatrix} \bmod 26 = \begin{bmatrix}21\\ 20\end{bmatrix} \equiv \textbf{V}$
\end{itemize}

\subsection{Contoh matriks non-invertible}

Sebagai ilustrasi syarat invertibilitas: matriks $K$ berukuran $2 \times 3$ (bukan persegi) tidak dapat memiliki invers. Untuk matriks $2 \times 2$ modulo 26, determinan harus relatif prima dengan 26. Contoh: $K=\begin{bmatrix}6 & 24\\ 13 & 16\end{bmatrix}$ mempunyai determinan $6 \cdot 16 - 24 \cdot 13 = 96 - 312 = -216$. Karena $\gcd(216, 26) = 2 \neq 1$, matriks tidak invertible modulo 26.

Cipher Hill dikembangkan oleh Lester S. Hill (1929) menggunakan aljabar linear dan modular arithmetic \cite{trappe2006}. Hill cipher rentan terhadap known-plaintext attack: dengan $n^2$ pasangan plaintext-ciphertext, penyerang dapat menyelesaikan sistem persamaan linear untuk memulihkan matriks kunci.

\subsection{Implementasi dalam C}

\begin{lstlisting}[style=cstyle, caption={Hill Cipher 2x2 dalam C}]
#include <stdio.h>

int mod_inverse(int a, int m) {
    int a0 = a, m0 = m, t = 0, new_t = 1;
    while (m0 != 0) {
        t = m0 / a0;
        m0 = a0 * t + m0 % a0;
        a0 = new_t;
        new_t = (a0 * t - m0 / a0) % a0;
    }
    return new_t < 0 ? new_t + m : new_t;
}

int matrix_inverse(int K[2][2], int invK[2][2], int n) {
    int det = K[0][0] * K[1][1] - K[0][1] * K[1][0];
    int det_inv = mod_inverse(det, n);
    int factor = det_inv;
    
    for (int i = 0; i < n; i++) {
        for (int j = 0; j < n; j++) {
            invK[i][j] = (factor * K[(j+1)%n][(i+1)%n]) % n;
        }
    }
    return 0;
}

void hill_encrypt(int block[2], int K[2][2], int cipher[2]) {
    cipher[0] = (K[0][0] * block[0] + K[0][1] * block[1]) % 26;
    cipher[1] = (K[1][0] * block[0] + K[1][1] * block[1]) % 26;
}

void hill_decrypt(int cipher[2], int invK[2][2], int plain[2]) {
    plain[0] = (invK[0][0] * cipher[0] + invK[0][1] * cipher[1]) % 26;
    plain[1] = (invK[1][0] * cipher[0] + invK[1][1] * cipher[1]) % 26;
}
\end{lstlisting}

Implementasi Hill cipher menunjukkan integrasi penuh antara konsep matematika dari Bab III dengan aplikasi kriptografi praktis. Fungsi \texttt{mod\_inverse} menggunakan Extended Euclidean Algorithm untuk menghitung invers modular, sedangkan \texttt{matrix\_inverse} mengimplementasikan adjoin kofaktor untuk invers matriks. Implementasi ini mendemonstrasikan bagaimana teori matriks dan aritmetika modular dapat dikombinasikan untuk menciptakan sistem enkripsi yang fungsional dan aman.

Dalam implementasi produksi, Hill cipher sering dioptimasi dengan menggunakan operasi matriks yang efisien dan representasi data yang terstruktur untuk meminimalkan overhead. Penting untuk dicatat bahwa implementasi Hill cipher yang benar harus memvalidasi parameter matriks kunci sebelum digunakan, karena matriks yang tidak invertible dapat menyebabkan kegagalan sistematis yang tidak dapat dipulihkan. Implementasi yang robust juga harus menangani berbagai ukuran blok dan kasus edge dengan baik.
